\chapter{Không Còn Một Khoảng Trống Nào}
\markboth{Không Còn Một Khoảng Trống Nào}{Không Còn Một Khoảng Trống Nào}

\begin{center}
	\textit{\color{bookprimary}``Nhiều người không nghiện thông tin. Họ nghiện cảm giác không phải ở một mình với đầu óc mình.''}

	\textit{\color{bookprimary}``Many people are not addicted to information. They are addicted to not being alone with their own minds.''}
\end{center}

\ngancach

\begin{tcolorbox}[colback=bookprimary!6,colframe=bookprimary!35,arc=5pt,title={\bfseries Góc Suy Nghĩ},fonttitle=\bfseries\color{bookprimary}]
Ngày xưa, con người có những khoảng trống tự nhiên.
\textit{(In the past, people had natural empty spaces.)}

Đứng chờ ai đó thì chờ.
\textit{(When they waited for someone, they simply waited.)}

Ngồi một mình thì ngồi.
\textit{(When they sat alone, they simply sat.)}

Đi vài bước thì đầu óc cũng đi cùng vài bước đó.
\textit{(When they walked a few steps, their minds walked with those steps too.)}

Bây giờ, rất nhiều học sinh và người lớn sống như thể mọi khoảng trống đều là một lỗi phải lấp ngay.
\textit{(Now many students and adults live as if every empty moment were an error that must be filled immediately.)}
\end{tcolorbox}

\section{Rút Điện Thoại Ra Như Một Phản Xạ}
\begin{truyen}{Rút Điện Thoại Ra Như Một Phản Xạ}{The Reflex of Reaching for the Screen}
\chuhoa{C}{ó một bạn đứng chờ trống vào lớp chưa tới ba mươi giây đã rút điện thoại ra.}
\textit{(A student waiting for the bell for less than thirty seconds already pulls out the phone.)}

Thầy hỏi: ``Em đang có việc gấp à?''
\textit{(The teacher asks: ``Do you have something urgent to do?'')}

Bạn ấy cười: ``Dạ không. Tại quen tay thôi thầy.''
\textit{(The student smiles: ``No. It is just a habit, teacher.'')}

\teacherpair{Đó mới là chỗ đáng ngại. Một thứ nguy hiểm nhất thường không còn cần lý do. Nó chỉ cần phản xạ. Khi em không còn chịu nổi ba mươi giây không chạm màn hình, vấn đề không còn là chiếc điện thoại nữa. Vấn đề là sức chịu yên của em đang teo lại rất nhanh.}{``That is exactly the worrying part. The most dangerous things are often the ones that no longer need reasons. They only need reflex. When you cannot endure thirty seconds without touching a screen, the problem is no longer just the phone. The problem is that your ability to stay still is shrinking very quickly.''}

Thầy nói tiếp: ``Người không chịu nổi khoảng trống rất dễ trở thành người không chịu nổi chiều sâu. Vì bất cứ việc gì cần nghĩ lâu hơn một chút, đau hơn một chút, chậm hơn một chút, họ cũng sẽ lập tức đi tìm một màn hình để trốn.''
\textit{(The teacher continues: ``A person who cannot endure emptiness often becomes a person who cannot endure depth. Anything that requires a little more thought, a little more pain, or a little more slowness will immediately send them running to a screen for escape.'')}
\end{truyen}

\section{Chờ Năm Phút Cũng Phải Lướt}
\begin{truyen}{Chờ Năm Phút Cũng Phải Lướt}{Needing to Scroll Through Every Five-Minute Wait}
\chuhoa{C}{ó bạn đứng đợi mẹ đón trước cổng trường, khoảng chờ chỉ vài phút nhưng mắt không rời nổi màn hình.}
\textit{(A student waits a few minutes at the school gate for a parent, yet cannot take their eyes off the screen.)}

Chờ đợi vốn là một phần bình thường của đời sống.
\textit{(Waiting is a normal part of life.)}

Nhưng nhiều người bây giờ coi nó như khoảng chết cần giết ngay.
\textit{(But many people now treat it like dead time that must be killed immediately.)}

\teacherpair{Em đừng coi vài phút chờ là vô ích. Chính những phút đó từng giúp con người quan sát, nhớ lại, nghĩ tiếp một điều dang dở. Nếu em giết hết những quãng đó bằng lướt, em sẽ dần mất luôn khả năng ở yên với một ý nghĩ chưa kịp chín.}{``Do not treat a few minutes of waiting as useless. Those minutes used to help people observe, remember, and continue a half-formed thought. If you kill all those moments with scrolling, you will gradually lose the ability to stay with an idea before it ripens.''}
\end{truyen}

\section{Vừa Xong Việc Là Mở Máy}
\begin{truyen}{Vừa Xong Việc Là Mở Máy}{Opening the Phone the Moment a Task Ends}
\chuhoa{C}{ó học sinh làm xong một bài tập nhỏ là lập tức chộp điện thoại, như thể đầu óc không được phép nghỉ nếu không có kích thích.}
\textit{(A student finishes one small task and immediately grabs the phone, as if the mind cannot rest without stimulation.)}

Khoảng chuyển giữa hai việc đáng lẽ giúp đầu óc đổi nhịp.
\textit{(The transition between two tasks should help the mind change rhythm.)}

Nhưng màn hình chen vào làm mọi nhịp nghỉ bị biến thành nhịp tiêu thụ.
\textit{(But the screen invades and turns every pause into consumption.)}

\teacherpair{Nghỉ không có nghĩa là nạp thêm tiếng ồn. Có những kiểu nghỉ thật sự chỉ xảy ra khi em để đầu óc thở một chút trước khi nhảy sang việc tiếp theo.}{``Rest does not mean adding more noise. Some forms of real rest happen only when you let the mind breathe a little before jumping to the next task.''}
\end{truyen}

\section{Xếp Hàng Không Chịu Được Im}
\begin{truyen}{Xếp Hàng Không Chịu Được Im}{Unable to Stay Quiet Even in a Queue}
\chuhoa{C}{ó bạn đứng xếp hàng mua nước chưa tới mười người mà cũng phải cúi xuống xem điện thoại từng nhịp.}
\textit{(A student stands in a short drink queue and still keeps checking the phone every few moments.)}

Thứ bị ăn mòn ở đây không chỉ là phép lịch sự.
\textit{(What erodes here is not only courtesy.)}

Nó còn là năng lực chịu đựng những nhịp chậm rất nhỏ của đời sống.
\textit{(It is also the capacity to endure life's very small slow rhythms.)}

\teacherpair{Người nào không chịu nổi một hàng đợi ngắn thì cũng khó chịu nổi một bài toán dài, một cuộc trò chuyện sâu, hay một nỗi buồn cần ngồi yên để đi qua.}{``A person who cannot endure a short queue will also struggle to endure a long problem, a deep conversation, or a sadness that needs stillness to pass.''}
\end{truyen}

\section{Đi Vài Bước Là Muốn Có Nhạc}
\begin{truyen}{Đi Vài Bước Là Muốn Có Nhạc}{Needing Sound for Every Few Steps}
\chuhoa{C}{ó học sinh đi từ lớp này sang lớp kia cũng phải cắm tai nghe cho đỡ trống.}
\textit{(A student moving from one class to another still needs earphones to avoid emptiness.)}

Âm thanh không xấu.
\textit{(Sound is not bad.)}

Nhưng nếu mọi quãng ngắn đều phải có thứ lấp vào, con người sẽ quên mất cảm giác nghe chính mình.
\textit{(But if every short interval must be filled, people forget how to hear themselves.)}

\teacherpair{Có những điều chỉ hiện ra khi xung quanh bớt ồn. Em không thể nghe được bên trong mình nếu suốt ngày tìm cách bị phủ kín từ bên ngoài.}{``Some things appear only when the surroundings quiet down. You cannot hear what is inside you if you spend all day trying to cover yourself from the outside.''}
\end{truyen}

\section{Trước Giờ Ngủ Vẫn Sợ Trống}
\begin{truyen}{Trước Giờ Ngủ Vẫn Sợ Trống}{Afraid of Emptiness Before Sleep}
\chuhoa{C}{ó bạn bảo rằng nếu không xem gì trước khi ngủ thì thấy rất khó chịu, như thể ngày chưa được khép lại.}
\textit{(A student says that without watching something before sleep, they feel uneasy, as if the day has not ended properly.)}

Thực ra đó không phải khép ngày.
\textit{(In truth, that is not closing the day.)}

Đó là kéo thêm kích thích vào đúng lúc đầu óc cần đi xuống.
\textit{(It is dragging in more stimulation at the exact time the mind should be descending.)}

\teacherpair{Người sợ vài phút trống trước khi ngủ thường không thiếu video. Họ thiếu khả năng tự ở cùng mình sau một ngày dài. Thiếu đó mới là thứ đáng lo.}{``A person who fears a few empty minutes before sleep is not lacking videos. They are lacking the ability to stay with themselves after a long day. That lack is the real concern.''}
\end{truyen}

\section{Mất Sóng Là Bồn Chồn}
\begin{truyen}{Mất Sóng Là Bồn Chồn}{Restless the Moment the Signal Disappears}
\chuhoa{C}{ó lúc mạng chập chờn vài phút mà cả nhóm học sinh đã bứt rứt như bị lấy mất một thứ rất lớn.}
\textit{(Sometimes the internet flickers for a few minutes and a whole group of students grows restless as if something huge has been taken away.)}

Sự bồn chồn đó nói nhiều hơn ta tưởng.
\textit{(That restlessness says more than it seems.)}

Nó cho thấy người ta đã quen dựa tinh thần vào một luồng kích thích luôn sẵn.
\textit{(It shows how used people have become to leaning emotionally on a constant stream of stimulation.)}

\teacherpair{Khi mất sóng mà em thấy hụt hẫng quá mạnh, có thể em không thiếu mạng. Em đang thiếu một đời sống bên trong đủ chắc để không sụp chỉ vì bên ngoài bỗng im đi.}{``If losing signal makes you feel too empty, you may not be lacking internet. You may be lacking an inner life sturdy enough not to collapse when the outside suddenly goes quiet.''}
\end{truyen}

\section{Ngồi Yên Ba Phút Cũng Khó}
\begin{truyen}{Ngồi Yên Ba Phút Cũng Khó}{Even Three Quiet Minutes Feel Hard}
\chuhoa{C}{ó thầy thử cho lớp ngồi yên ba phút trước giờ học, không điện thoại, không nói chuyện, chỉ thở và nhìn lại mình.}
\textit{(A teacher asks the class to sit quietly for three minutes before the lesson, with no phones and no talking, only breathing and returning to themselves.)}

Nhiều em cười trừ, xoay bút, nhúc nhích, nhìn đồng hồ như vừa bị phạt rất nặng.
\textit{(Many students laugh awkwardly, spin pens, fidget, and stare at the clock as if severely punished.)}

\teacherpair{Nếu ba phút yên cũng khó chịu, em đừng coi thường chuyện đó. Một đầu óc không ở yên nổi rất khó học sâu, khó cầu nguyện, khó suy nghĩ, và cũng khó trưởng thành tử tế.}{``If even three quiet minutes feel unbearable, do not dismiss that. A mind that cannot stay still will struggle to learn deeply, pray, reflect, and mature decently.''}
\end{truyen}

\begin{baihoc}
Điện thoại thắng con người không phải ở những giờ rất dài.
\textit{(The phone defeats people not through very long hours alone.)}

Nó thắng ở những khoảng cực ngắn mà con người không còn chịu nổi sự không được kích thích.
\textit{(It wins in very short gaps when people can no longer endure the absence of stimulation.)}
\end{baihoc}

\begin{ghinhoanh}
\textbf{Short English to remember:}

Silence is not emptiness.

It is where your mind learns to stand.
\end{ghinhoanh}

\begin{tuhocnhanh}
1. Em có thể chờ năm phút mà không mở điện thoại không? \textit{(Can you wait five minutes without opening your phone?)}

2. Khoảng trống nào trong ngày em sợ nhất? \textit{(Which empty moment of the day do you fear most?)}

3. Em đang dùng màn hình để làm việc, hay để khỏi phải ở cùng đầu óc mình? \textit{(Are you using the screen to work, or to avoid being with your own mind?)}
\end{tuhocnhanh}

\begin{tongketchuong}
Người không chịu nổi khoảng trống rất dễ trở thành người không chịu nổi chiều sâu.
\textit{(A person who cannot bear emptiness easily becomes a person who cannot bear depth.)}
\end{tongketchuong}

\ngancach